\documentclass[10pt]{article}
\usepackage{graphicx} % Required for inserting images
\usepackage{hyperref} % To include links
\usepackage{anyfontsize}
\usepackage[a4paper, margin = 2.5cm]{geometry}
\usepackage{mwe}
\usepackage{fancyhdr}
\usepackage{amsmath}
\usepackage{setspace}
\usepackage{indentfirst}
\renewcommand{\baselinestretch}{1.25}
\usepackage{xcolor}
\definecolor{durisimo}{RGB}{0,160,255}
\definecolor{durisimo2}{RGB}{7,114,183}
\definecolor{durisimo3}{RGB}{0,180,255}
\definecolor{durisimo4}{RGB}{0, 210, 255}
\usepackage{tikz}
\usepackage[nottoc]{tocbibind}


\pagenumbering{Roman}
\title{image\_project}
\author{Héctor }
\date{April 2026}

\begin{document}
    
    %TITLE STARTS
    \begin{titlepage}
        \thispagestyle{empty}
        \newgeometry{left=2cm, right=1cm, top=2cm, bottom=2cm}
        \begin{tikzpicture}[overlay]
            \fill[durisimo2] (-10,-3) rectangle (19,-11.05)
            \fill[durisimo] (-10, -11) rectangle (19, -11.91)
            \fill[durisimo3] (-10, -11.9) rectangle (19, -12.751)
            \fill[durisimo4] (-10, -12.75) rectangle (19, -13.25)

            \fill[durisimo2] (-10, -20) rectangle (19, -28)
        \end{tikzpicture}
        \begin{spacing}{1}
            \begin{flushleft}
            \vspace*{3cm}
            \fontsize{55}{65}\selectfont {\rmfamily \bfseries OPTICAL \newline CHARACTER\newline RECOGNITION\par}
            \end{flushleft}
        \end{spacing}
        \vspace{3cm}
        \begin{spacing}{1}
            \begin{flushleft}
                {\huge {Fundamentals of Sound and Image project 2025-2026}}
            \end{flushleft}
        \end{spacing}
        \vspace{6cm}
        \begin{flushleft}
        \begin{spacing}{1.5}
            {\huge \textsc{Jorge López Cuns \newline
            Pablo Sanles Santos \newline
            Hugo Mera Tato \newline
            Héctor Costa Portas}}
        \end{spacing}
        \end{flushleft}
        \begin{flushright}
            \Large{\textbf{Universidade de Vigo}}
        \end{flushright}
    \end{titlepage}
    %TITLE ENDS
    %DEFINE THE STYLE FOR REGULAR PAGES
    \pagestyle{fancy}
        \fancyhf{}
        \fancyfoot[R]{\thepage}
        \fancyfoot[L]{Image project 2025-2026. University of Vigo}
        \fancyhead[C]

\renewcommand*\contentsname{Table of contents}
\tableofcontents
\newpage
\section{Introduction}
\section{{Task 1: Binarization}}
    \indent Our first task is related with the conversion of the given image to a binary format, that is, black and white. We will be defining a method to determine whether a given pixel will be converted to black or to white depending on a few parameters. \newline \newline 
    \begin{enumerate}
        \item First of all, we convert the image from RGB to grayscale via $rgb2gray.m$, a function already provided by MATLAB.
        \item Then, we must define a standard to decide if we want to convert a given pixel to white or to black. In other words, a \textbf{threshold}. For example, we can take the mean of the luminance of all pixels in the image. However, there is a function in MATLAB that has already optimized this process, called $graythresh.m$, which is based on the MATLAB function $otsuthresh.m$. \small{\cite{otsu}}
        \item Once we have a grayscale image and a threshold to have a black-white-defining criterium, we can use $imbinarize.m$ to perform the binarization with these two parameters.
    \end{enumerate} 


\newpage

\begin{thebibliography}

    \bibitem{otsu}
        \textit{Otsu thresholding:}  \newline \textcolor{blue}{\url{https://engineering.purdue.edu/kak/computervision/ECE661.08/OTSU_paper.pdf}}
    \addcontentsline{toc}{section}{References}

\end{thebibliography}
\end{document}
